\chapter{Phase 1: Lokaler Kern}

\section{Ziel}

Ein lokales Repository mit einem Nutzer. Es lässt sich anlegen, Dateien werden gespeichert, Commits erzeugt und der Status angezeigt.

\section{Umfang}

\startitemize
\item Objektspeicher mit Blobs, Trees und Commits, adressiert über BLAKE3 mit Präfix nach ADR-001.
\item Referenzen für den Branch \type{main} und HEAD.
\item Index als Staging-Bereich.
\item Befehle \type{init}, \type{add}, \type{commit} und \type{status}.
\item Akteur-Felder in jedem Commit und im Operation Log nach ADR-002.
\stopitemize

\section{Aufbau eines Repositories}

Ein Repository ist ein Verzeichnis, das einen Ordner \type{.stemma} enthält. Die Schicht \type{Repository} in \type{stemma-core} bündelt Objektspeicher und Referenzen, damit jede Operation die Speicher desselben Repositories verwendet (ADR-009).

\starttyping
.stemma/
    format        stemma-repository 1
    HEAD          branch main
    objects/      lose Objekte
    refs/heads/   Branches
    tmp/          temporäre Dateien des Objektspeichers
\stoptyping

Die Datei \type{format} trägt eine Versionsnummer. Beim Öffnen lehnt der Kern eine fehlende, eine nicht kanonische und jede andere Version ab. Ein älteres Programm weist damit ein neueres Repository zurück, statt es zu beschädigen.

Der Befehl \type{init} baut den Ordner unter einem temporären Namen neben \type{.stemma} auf und benennt ihn am Ende um. Ein Abbruch hinterlässt so kein halbes Repository. Existiert \type{.stemma} bereits, bricht \type{init} ab und ändert nichts. Der Standard-Branch ist \type{main} und lässt sich mit \type{--initial-branch} wählen. HEAD zeigt auf ihn, der Branch selbst entsteht erst mit dem ersten Commit. Wie jede Operation ruft \type{init} zuerst \type{authorize} mit dem Akteur auf.

\section{Offene Fragen}

\startitemize
\item Format des Index: eigenes Binärformat oder lesbares Textformat für Debugging?
\item Ausführbare Dateien und Dateirechte: Welche Metadaten werden gespeichert?
\item Symbolische Links: Speichern oder ablehnen?
\item Arbeitskopie als Commit, wie bei Jujutsu, oder klassischer Index?
\stopitemize
